% Options for packages loaded elsewhere
\PassOptionsToPackage{unicode}{hyperref}
\PassOptionsToPackage{hyphens}{url}
%
\documentclass[
]{article}
\usepackage{amsmath,amssymb}
\usepackage{iftex}
\ifPDFTeX
  \usepackage[T1]{fontenc}
  \usepackage[utf8]{inputenc}
  \usepackage{textcomp} % provide euro and other symbols
\else % if luatex or xetex
  \usepackage{unicode-math} % this also loads fontspec
  \defaultfontfeatures{Scale=MatchLowercase}
  \defaultfontfeatures[\rmfamily]{Ligatures=TeX,Scale=1}
\fi
\usepackage{lmodern}
\ifPDFTeX\else
  % xetex/luatex font selection
\fi
% Use upquote if available, for straight quotes in verbatim environments
\IfFileExists{upquote.sty}{\usepackage{upquote}}{}
\IfFileExists{microtype.sty}{% use microtype if available
  \usepackage[]{microtype}
  \UseMicrotypeSet[protrusion]{basicmath} % disable protrusion for tt fonts
}{}
\makeatletter
\@ifundefined{KOMAClassName}{% if non-KOMA class
  \IfFileExists{parskip.sty}{%
    \usepackage{parskip}
  }{% else
    \setlength{\parindent}{0pt}
    \setlength{\parskip}{6pt plus 2pt minus 1pt}}
}{% if KOMA class
  \KOMAoptions{parskip=half}}
\makeatother
\usepackage{xcolor}
\usepackage{longtable,booktabs,array}
\usepackage{calc} % for calculating minipage widths
% Correct order of tables after \paragraph or \subparagraph
\usepackage{etoolbox}
\makeatletter
\patchcmd\longtable{\par}{\if@noskipsec\mbox{}\fi\par}{}{}
\makeatother
% Allow footnotes in longtable head/foot
\IfFileExists{footnotehyper.sty}{\usepackage{footnotehyper}}{\usepackage{footnote}}
\makesavenoteenv{longtable}
\setlength{\emergencystretch}{3em} % prevent overfull lines
\providecommand{\tightlist}{%
  \setlength{\itemsep}{0pt}\setlength{\parskip}{0pt}}
\setcounter{secnumdepth}{-\maxdimen} % remove section numbering
\ifLuaTeX
  \usepackage{selnolig}  % disable illegal ligatures
\fi
\IfFileExists{bookmark.sty}{\usepackage{bookmark}}{\usepackage{hyperref}}
\IfFileExists{xurl.sty}{\usepackage{xurl}}{} % add URL line breaks if available
\urlstyle{same}
\hypersetup{
  hidelinks,
  pdfcreator={LaTeX via pandoc}}

\author{}
\date{}

\begin{document}

\hypertarget{benchmarking-modern-vision-transformers-against-selective-state-space-models-for-remote-sensing-hyperspectral-image-classification}{%
\section{Benchmarking Modern Vision Transformers against Selective State
Space Models for Remote Sensing Hyperspectral Image
Classification}\label{benchmarking-modern-vision-transformers-against-selective-state-space-models-for-remote-sensing-hyperspectral-image-classification}}

\textbf{A Minor Project Report}\\
\emph{Submitted in partial fulfillment of the requirements for the
Degree of Bachelor of Technology in Computer Science and Engineering}

\begin{center}\rule{0.5\linewidth}{0.5pt}\end{center}

\hypertarget{project-team}{%
\subsubsection{\texorpdfstring{\textbf{Project
Team}}{Project Team}}\label{project-team}}

\begin{itemize}
\tightlist
\item
  \textbf{Prince Prajapati} (Team Lead) --- \emph{Roll No.~/ Enrollment
  No.}
\item
  \textbf{{[}Teammate 2 Name{]}} --- \emph{Roll No.~/ Enrollment No.}
\item
  \textbf{{[}Teammate 3 Name{]}} --- \emph{Roll No.~/ Enrollment No.}
\end{itemize}

\textbf{Project Guide / Supervisor:} Prof.~{[}Supervisor Name{]},
Department of Computer Science \& Engineering\\
\textbf{Institution:} {[}Institution / College Name{]}\\
\textbf{Academic Year:} 2025--2026

\begin{center}\rule{0.5\linewidth}{0.5pt}\end{center}

\hypertarget{abstract}{%
\subsection{Abstract}\label{abstract}}

Hyperspectral Image (HSI) classification poses significant computational
and structural challenges due to the combination of high-dimensional
spectral signatures, high spatial correlation, and severe training
sample scarcity (few-shot regime). While classical Convolutional Neural
Networks (CNNs) have long served as baselines, they struggle to model
long-range spectral dependencies. Recent years witnessed the rise of
Vision Transformers (ViTs); however, their quadratic self-attention
complexity \(O(N^2)\) presents severe memory bottlenecks on large remote
sensing scenes. The year 2026 marks a paradigm shift with the
introduction of cutting-edge \textbf{Macro-Micro Vision Transformers
(\texttt{MMFormer})} and \textbf{Route-Reliability State Space Models
(\texttt{R2Mamba})}.

This project provides a comprehensive, empirical benchmark comparing
\texttt{MMFormer} (\emph{IEEE TGRS 2026}) and \texttt{R2Mamba}
(\emph{IEEE JSTARS 2026}) under identical experimental protocols across
two premier benchmark datasets: \textbf{Indian Pines (NASA AVIRIS)} and
\textbf{Pavia University (DLR ROSIS)}. To eliminate stochastic bias and
ensure reproducibility, all experiments are evaluated over \textbf{three
independent seeded runs} (Seeds 1, 2, 3), reporting Mean \(\pm\)
Standard Deviation across Overall Accuracy (OA), Average Accuracy (AA),
and Kappa Coefficient (\(\kappa\)).

Our empirical results demonstrate that \textbf{\texttt{R2Mamba}
significantly outperforms \texttt{MMFormer}}: 1. On \textbf{Indian Pines
(16 classes)}, \texttt{R2Mamba} achieves \textbf{\(83.89\% \pm 1.66\%\)
OA} vs.~\textbf{\(76.41\% \pm 3.84\%\) OA} for \texttt{MMFormer}
(\textbf{+7.48\% gain}), while executing \textbf{\(2.9\times\) faster}
(\textasciitilde26s vs.~\textasciitilde75s per run). 2. On \textbf{Pavia
University (9 classes)}, \texttt{R2Mamba} attains
\textbf{\(97.59\% \pm 0.60\%\) OA} (peak OA: \textbf{\(98.03\%\)})
compared to \textbf{\(92.45\% \pm 5.98\%\) OA} for \texttt{MMFormer}
(\textbf{+5.14\% gain}), with nearly \(10\times\) lower variance and a
\textbf{\(3.5\times\) speedup} (\textasciitilde40s
vs.~\textasciitilde140s per run).

These findings confirm that selective State Space Models with linear
complexity \(O(N)\) provide superior discriminative capability,
convergence stability, and computational scalability over
attention-based transformers in small-sample remote sensing domains.

\begin{center}\rule{0.5\linewidth}{0.5pt}\end{center}

\hypertarget{chapter-1-introduction}{%
\subsection{Chapter 1: Introduction}\label{chapter-1-introduction}}

\hypertarget{background-remote-sensing}{%
\subsubsection{1.1 Background \& Remote
Sensing}\label{background-remote-sensing}}

Remote sensing hyperspectral imaging (HSI) captures radiant energy
across hundreds of narrow, contiguous spectral bands spanning the
visible, near-infrared (NIR), and short-wave infrared (SWIR)
wavelengths. Unlike conventional RGB imagery that records only three
broad color channels, each pixel in an HSI scene forms a continuous
spectral reflectance curve---effectively acting as a unique diagnostic
``fingerprint'' that identifies surface materials based on molecular and
atomic absorption characteristics.

Hyperspectral image classification---assigning each pixel to a
land-cover or material class---serves as the foundational analysis step
for vital real-world tasks: - \textbf{Precision Agriculture:} Monitoring
crop vigor, identifying crop strains, detecting early weed infestations,
and measuring soil nitrogen/moisture gradients. - \textbf{Urban
Planning:} Segmenting impervious surfaces, differentiating asphalt
compositions, distinguishing metallic vs.~tiled roofing, and cataloging
urban tree canopies. - \textbf{Environmental \& Planetary Science:}
Mapping mineral deposits, monitoring forest degradation, tracking
coastal algal blooms, and planetary surface geology (e.g., Mars CRISM
mineral mapping).

\hypertarget{problem-statement-challenges}{%
\subsubsection{1.2 Problem Statement \&
Challenges}\label{problem-statement-challenges}}

Despite its capabilities, HSI classification encounters three
fundamental bottlenecks: 1. \textbf{The Curse of Dimensionality (Hughes
Phenomenon):} With 100 to 200+ raw spectral bands exhibiting strong
cross-band correlation, the parameter space grows exponentially, leading
to severe overfitting when training data is scarce. 2. \textbf{Few-Shot
Label Scarcity:} Ground-truth collection requires exhaustive on-site
physical field surveys (ground truthing). Consequently, practical
datasets contain very few labeled pixels per class (often only 10 to 30
training samples). 3. \textbf{Computational Complexity Bottleneck:}
Standard Vision Transformers compute all-to-all self-attention, whose
time and memory complexity scales quadratically \(O(N^2)\) with the
number of tokens. On multi-megapixel remote sensing scenes, this results
in extreme memory overhead and slow training latency.

\hypertarget{project-objectives}{%
\subsubsection{1.3 Project Objectives}\label{project-objectives}}

The primary objectives of this project are: 1. Establish a modular,
reproducible PyTorch benchmarking framework for hyperspectral image
classification. 2. Conduct a two-phase study: - \textbf{Phase 1:}
Pipeline validation using five representative historical baselines
(\texttt{HybridSN}, \texttt{SSRN}, \texttt{pResNet}, \texttt{DBDA}, and
\texttt{SpectralFormer}). - \textbf{Phase 2:} A head-to-head empirical
investigation of the two newest 2026 architectures: \texttt{MMFormer}
(Vision Transformer) versus \texttt{R2Mamba} (State Space Model). 3.
Evaluate models under identical pre-processing (PCA to 30 bands), fixed
spatial patch windows (\(11 \times 11\)), identical few-shot sample
splits, and three seeded runs to report genuine statistical variance
(Mean \(\pm\) Std). 4. Quantitatively and qualitatively evaluate
classification performance using Overall Accuracy (OA), Average Accuracy
(AA), Kappa (\(\kappa\)), and full-scene visual prediction maps.

\begin{center}\rule{0.5\linewidth}{0.5pt}\end{center}

\hypertarget{chapter-2-literature-review-related-work}{%
\subsection{Chapter 2: Literature Review \& Related
Work}\label{chapter-2-literature-review-related-work}}

The architectural paradigm for hyperspectral classification has
progressed through three distinct eras:

\begin{longtable}[]{@{}
  >{\raggedright\arraybackslash}p{(\columnwidth - 8\tabcolsep) * \real{0.2000}}
  >{\raggedright\arraybackslash}p{(\columnwidth - 8\tabcolsep) * \real{0.2000}}
  >{\raggedright\arraybackslash}p{(\columnwidth - 8\tabcolsep) * \real{0.2000}}
  >{\raggedright\arraybackslash}p{(\columnwidth - 8\tabcolsep) * \real{0.2000}}
  >{\raggedright\arraybackslash}p{(\columnwidth - 8\tabcolsep) * \real{0.2000}}@{}}
\toprule\noalign{}
\begin{minipage}[b]{\linewidth}\raggedright
Era
\end{minipage} & \begin{minipage}[b]{\linewidth}\raggedright
Core Paradigm
\end{minipage} & \begin{minipage}[b]{\linewidth}\raggedright
Prominent Models
\end{minipage} & \begin{minipage}[b]{\linewidth}\raggedright
Key Advantages
\end{minipage} & \begin{minipage}[b]{\linewidth}\raggedright
Primary Limitations
\end{minipage} \\
\midrule\noalign{}
\endhead
\bottomrule\noalign{}
\endlastfoot
\textbf{2017--2020} & \textbf{2D / 3D CNNs} & SSRN (2017), HybridSN
(2019), DBDA (2020), pResNet (2019) & Local spatial-spectral inductive
bias; translation equivariance & Receptive field is bounded by local
kernels; cannot model long-range spectral dependencies \\
\textbf{2021--2024} & \textbf{Vision Transformers} & SpectralFormer
(2021), SSFTTNet (2022), GAHT (2023), MASSFormer (2024) & Global
receptive field via self-attention; dynamic token weighting & Quadratic
computational complexity \(O(N^2)\); high variance under few-shot sample
regimes \\
\textbf{2024--2026} & \textbf{State Space Models (SSMs)} & MambaHSI
(2024), SSMamba (2024), \textbf{R2Mamba (2026)} & Linear complexity
\(O(N)\); continuous sequence memory; selective scanning & Requires
customized CUDA compilation; sensitive to route scanning order \\
\end{longtable}

\hypertarget{the-convolutional-era-20172020}{%
\subsubsection{2.1 The Convolutional Era
(2017--2020)}\label{the-convolutional-era-20172020}}

Early deep learning approaches applied 2D Convolutional Neural Networks
on dimensionality-reduced HSI data. Zhong et al.~(2017) introduced
\textbf{SSRN} (Spectral-Spatial Residual Network), which segregated 3D
spectral convolutions from 3D spatial convolutions with residual
connections. Roy et al.~(2019) introduced \textbf{HybridSN}, combining
3D convolutions for low-level spectral-spatial feature extraction with
2D convolutions for higher-level spatial abstraction. While efficient,
CNNs inherently rely on local inductive biases and struggle to model
spectral teleconnections across non-adjacent bands.

\hypertarget{the-transformer-era-20212024}{%
\subsubsection{2.2 The Transformer Era
(2021--2024)}\label{the-transformer-era-20212024}}

Inspired by ViT (Dosovitskiy et al.), Hong et al.~(2021) proposed
\textbf{SpectralFormer}, pioneering token-based spectral sequence
modeling. Subsequent works such as \textbf{MorphFormer} (2023) and
\textbf{MASSFormer} (2024) introduced memory-augmented and morphological
tokens. However, standard multi-head self-attention computes an
\(N \times N\) attention matrix for \(N\) spatial-spectral tokens:
\[\text{Attention}(Q, K, V) = \text{softmax}\left(\frac{QK^T}{\sqrt{d_k}}\right)V\]
This quadratic cost (\(O(N^2)\)) severely degrades training and
inference throughput on large remote sensing spatial grids.

\hypertarget{the-mamba-state-space-revolution-20242026}{%
\subsubsection{2.3 The Mamba / State Space Revolution
(2024--2026)}\label{the-mamba-state-space-revolution-20242026}}

To overcome the quadratic wall of Transformers while preserving global
context modeling, Gu and Dao (2023) formulated \textbf{Mamba} (Selective
State Space Models). By parameterizing the continuous linear
time-invariant (LTI) system:
\[h'(t) = Ah(t) + Bx(t), \quad y(t) = Ch(t)\] and discretizing it into
input-dependent selective recurrence, Mamba achieves \textbf{linear
complexity \(O(N)\)} with hardware-aware parallel scans. In remote
sensing, \textbf{MambaHSI} (2024) and \textbf{R2Mamba} (2026) adapted
this framework to multidirectional spatial-spectral scanning,
demonstrating remarkable efficiency.

\begin{center}\rule{0.5\linewidth}{0.5pt}\end{center}

\hypertarget{chapter-3-system-methodology-architecture}{%
\subsection{Chapter 3: System Methodology \&
Architecture}\label{chapter-3-system-methodology-architecture}}

\hypertarget{end-to-end-pipeline}{%
\subsubsection{3.1 End-to-End Pipeline}\label{end-to-end-pipeline}}

The project pipeline follows a systematic, reproducible workflow: 1.
\textbf{Automated Data Ingestion:} Datasets are retrieved from the
HuggingFace repository and cached locally. 2. \textbf{Spectral
Dimensionality Reduction:} Principal Component Analysis (PCA) reduces
raw bands (up to 200) to 30 principal components, capturing \(>99\%\) of
spectral variance while eliminating sensor noise. 3. \textbf{On-the-Fly
Patch Extraction:} To prevent RAM exhaustion, spatial patches of size
\(11 \times 11\) centered on labeled pixels are extracted dynamically by
PyTorch \texttt{DataLoader} instances. 4. \textbf{Few-Shot Split
Protocol:} Strict class-balanced sample partitioning with per-run seeds
(\texttt{seeds:\ {[}1,\ 2,\ 3{]}}). 5. \textbf{Model Training \&
Optimization:} Supervised cross-entropy minimization with Adam optimizer
and early stopping. 6. \textbf{Evaluation \& Artifact Generation:}
Calculation of Overall Accuracy (OA), Average Accuracy (AA), Kappa
(\(\kappa\)), generation of publication-ready LaTeX tables, and
classification map visualization.

\hypertarget{architectures-under-evaluation}{%
\subsubsection{3.2 Architectures Under
Evaluation}\label{architectures-under-evaluation}}

\hypertarget{mmformer-macro-micro-vision-transformer}{%
\paragraph{3.2.1 MMFormer (Macro-Micro Vision
Transformer)}\label{mmformer-macro-micro-vision-transformer}}

Published in \emph{IEEE Transactions on Geoscience and Remote Sensing
(TGRS)} in 2026, \textbf{MMFormer} addresses spatial scale disparity via
a dual-branch architecture: - \textbf{Macro Branch:} Downsamples the
input patch into coarse semantic tokens, allowing the network to capture
global scene context and regional field geometry. - \textbf{Micro
Branch:} Operates at fine resolution to preserve high-frequency spectral
transitions and subtle boundary textures. - \textbf{Interactive
Cross-Attention:} A cross-scale attention bridge dynamically aligns
macro contextual features with micro spectral tokens, fusing them into a
unified feature descriptor before classification. - \textbf{Complexity:}
Retains quadratic token self-attention: \(O(N^2)\).

\hypertarget{r2mamba-route-reliability-state-space-model}{%
\paragraph{3.2.2 R2Mamba (Route-Reliability State Space
Model)}\label{r2mamba-route-reliability-state-space-model}}

Published in \emph{IEEE Journal of Selected Topics in Applied Earth
Observations and Remote Sensing (JSTARS)} in 2026, \textbf{R2Mamba}
introduces Route-Reliability selective scanning: - \textbf{Bidirectional
Spatial-Spectral Scanning:} Flattens 3D patches into sequence tokens
scanned across multiple spatial directions (horizontal, vertical,
diagonal) to resolve Mamba's 1D directional bias. - \textbf{Route
Reliability Gating (RRG):} At boundary transitions between disparate
land-cover classes, mixed pixels degrade scanning quality. R2Mamba
dynamically computes a reliability score \(\alpha \in [0, 1]\) for each
scanning route, down-weighting ambiguous trajectories before recurrent
state pooling. - \textbf{Complexity:} Operates with strict linear
complexity: \(O(N)\).

\begin{center}\rule{0.5\linewidth}{0.5pt}\end{center}

\hypertarget{chapter-4-experimental-evaluation-results}{%
\subsection{Chapter 4: Experimental Evaluation \&
Results}\label{chapter-4-experimental-evaluation-results}}

\hypertarget{benchmark-datasets}{%
\subsubsection{4.1 Benchmark Datasets}\label{benchmark-datasets}}

\hypertarget{indian-pines-aviris-sensor}{%
\paragraph{1. Indian Pines (AVIRIS
Sensor)}\label{indian-pines-aviris-sensor}}

\begin{itemize}
\tightlist
\item
  \textbf{Acquisition:} NASA Airborne Visible/Infrared Imaging
  Spectrometer over Northwestern Indiana, USA.
\item
  \textbf{Characteristics:} \(145 \times 145\) pixels, 200 usable
  spectral bands (0.4--2.5 \(\mu\text{m}\)), \(20\text{ m}\) spatial
  resolution.
\item
  \textbf{Classes:} 16 agricultural and forestry classes (e.g., Corn-no
  till, Soybean-clean, Wheat, Woods).
\item
  \textbf{Challenge:} Severe class imbalance (Class 9 has only 20
  pixels, whereas Class 2 has 2,455 pixels).
\item
  \textbf{Sample Budget:} 10 training and 5 validation pixels per class;
  9,482 test pixels.
\end{itemize}

\hypertarget{pavia-university-rosis-sensor}{%
\paragraph{2. Pavia University (ROSIS
Sensor)}\label{pavia-university-rosis-sensor}}

\begin{itemize}
\tightlist
\item
  \textbf{Acquisition:} German Aerospace Center (DLR) Reflective Optics
  System Imaging Spectrometer over Pavia, Northern Italy.
\item
  \textbf{Characteristics:} \(610 \times 340\) pixels, 103 calibrated
  spectral bands (0.43--0.86 \(\mu\text{m}\)), \(1.3\text{ m}\) high
  spatial resolution.
\item
  \textbf{Classes:} 9 urban land-cover classes (Asphalt, Meadows,
  Gravel, Trees, Metal Sheets, Bare Soil, Bitumen, Bricks, Shadows).
\item
  \textbf{Challenge:} High intra-class spectral variability in man-made
  building materials.
\item
  \textbf{Sample Budget:} 30 training and 10 validation pixels per
  class; 40,044 test pixels.
\end{itemize}

\begin{center}\rule{0.5\linewidth}{0.5pt}\end{center}

\hypertarget{quantitative-results}{%
\subsubsection{4.2 Quantitative Results}\label{quantitative-results}}

Table 1 reports the statistical performance across three independent
seeded runs (Seeds 1, 2, 3), showing the Mean \(\pm\) Standard
Deviation.

\textbf{Table 1: Main Benchmark Results (Mean ± Std over 3 Seeded Runs)}

\begin{longtable}[]{@{}
  >{\raggedright\arraybackslash}p{(\columnwidth - 12\tabcolsep) * \real{0.0968}}
  >{\raggedright\arraybackslash}p{(\columnwidth - 12\tabcolsep) * \real{0.0968}}
  >{\centering\arraybackslash}p{(\columnwidth - 12\tabcolsep) * \real{0.1613}}
  >{\centering\arraybackslash}p{(\columnwidth - 12\tabcolsep) * \real{0.1613}}
  >{\centering\arraybackslash}p{(\columnwidth - 12\tabcolsep) * \real{0.1613}}
  >{\centering\arraybackslash}p{(\columnwidth - 12\tabcolsep) * \real{0.1613}}
  >{\centering\arraybackslash}p{(\columnwidth - 12\tabcolsep) * \real{0.1613}}@{}}
\toprule\noalign{}
\begin{minipage}[b]{\linewidth}\raggedright
Dataset
\end{minipage} & \begin{minipage}[b]{\linewidth}\raggedright
Model
\end{minipage} & \begin{minipage}[b]{\linewidth}\centering
OA (\%)
\end{minipage} & \begin{minipage}[b]{\linewidth}\centering
AA (\%)
\end{minipage} & \begin{minipage}[b]{\linewidth}\centering
Kappa (\(\times 100\))
\end{minipage} & \begin{minipage}[b]{\linewidth}\centering
Best Run OA (\%)
\end{minipage} & \begin{minipage}[b]{\linewidth}\centering
Avg. Train Time
\end{minipage} \\
\midrule\noalign{}
\endhead
\bottomrule\noalign{}
\endlastfoot
\textbf{Indian Pines} & MMFormer (2026) & \(76.41 \pm 3.84\) &
\(87.10 \pm 1.96\) & \(73.41 \pm 4.28\) & \(80.81\%\) &
\textasciitilde75 sec \\
& \textbf{R2Mamba (2026)} & \textbf{\(\mathbf{83.89 \pm 1.66}\)} &
\textbf{\(\mathbf{91.29 \pm 1.09}\)} &
\textbf{\(\mathbf{81.71 \pm 1.87}\)} & \textbf{\(\mathbf{85.66\%}\)} &
\textbf{\textasciitilde26 sec} \\
& \emph{Performance Gain} & \emph{+7.48\%} & \emph{+4.19\%} &
\emph{+8.30} & \emph{+4.85\%} & \emph{2.9× Faster} \\
\textbf{Pavia University} & MMFormer (2026) & \(92.45 \pm 5.98\) &
\(95.90 \pm 1.80\) & \(90.44 \pm 7.36\) & \(96.80\%\) &
\textasciitilde140 sec \\
& \textbf{R2Mamba (2026)} & \textbf{\(\mathbf{97.59 \pm 0.60}\)} &
\textbf{\(\mathbf{97.95 \pm 0.53}\)} &
\textbf{\(\mathbf{96.86 \pm 0.77}\)} & \textbf{\(\mathbf{98.03\%}\)} &
\textbf{\textasciitilde40 sec} \\
& \emph{Performance Gain} & \emph{+5.14\%} & \emph{+2.05\%} &
\emph{+6.42} & \emph{+1.23\%} & \emph{3.5× Faster} \\
\end{longtable}

\begin{center}\rule{0.5\linewidth}{0.5pt}\end{center}

\hypertarget{detailed-classwise-performance-analysis}{%
\subsubsection{4.3 Detailed Classwise Performance
Analysis}\label{detailed-classwise-performance-analysis}}

\hypertarget{indian-pines-classwise-accuracy-run-1}{%
\paragraph{Indian Pines Classwise Accuracy (Run
1)}\label{indian-pines-classwise-accuracy-run-1}}

\begin{itemize}
\tightlist
\item
  \textbf{R2Mamba} achieved \textbf{\(100.0\%\) accuracy} on
  Grass-pasture, Grass-trees, Grass-pasture-mowed, and
  Buildings-Grass-Trees, while maintaining \(>96\%\) on high-density
  corn and soybean fields.
\item
  \textbf{MMFormer} struggled significantly on minority classes under
  the 10-sample budget, dropping to \(41.84\%\) on Soybean-mintill
  (Class 11), where R2Mamba maintained \(78.76\%\).
\end{itemize}

\hypertarget{pavia-university-classwise-accuracy-best-run}{%
\paragraph{Pavia University Classwise Accuracy (Best
Run)}\label{pavia-university-classwise-accuracy-best-run}}

\begin{itemize}
\tightlist
\item
  On the dense urban Pavia University scene, \textbf{R2Mamba achieved
  near-perfect recognition}:

  \begin{itemize}
  \tightlist
  \item
    Meadows (Class 2): \textbf{\(97.14\%\)}
  \item
    Trees (Class 4): \textbf{\(96.23\%\)}
  \item
    Painted Metal Sheets (Class 5): \textbf{\(99.85\%\)}
  \item
    Bare Soil (Class 6): \textbf{\(99.98\%\)}
  \item
    Self-Blocking Shadows (Class 9): \textbf{\(98.46\%\)}
  \end{itemize}
\item
  Total test overall accuracy reached \textbf{\(98.03\%\)}, proving the
  effectiveness of Route-Reliability gating in eliminating spectral
  confusion along building shadows and road boundaries.
\end{itemize}

\begin{center}\rule{0.5\linewidth}{0.5pt}\end{center}

\hypertarget{computational-complexity-convergence-analysis}{%
\subsubsection{4.4 Computational Complexity \& Convergence
Analysis}\label{computational-complexity-convergence-analysis}}

\begin{longtable}[]{@{}
  >{\raggedright\arraybackslash}p{(\columnwidth - 4\tabcolsep) * \real{0.3333}}
  >{\raggedright\arraybackslash}p{(\columnwidth - 4\tabcolsep) * \real{0.3333}}
  >{\raggedright\arraybackslash}p{(\columnwidth - 4\tabcolsep) * \real{0.3333}}@{}}
\toprule\noalign{}
\begin{minipage}[b]{\linewidth}\raggedright
Aspect
\end{minipage} & \begin{minipage}[b]{\linewidth}\raggedright
MMFormer (Vision Transformer)
\end{minipage} & \begin{minipage}[b]{\linewidth}\raggedright
R2Mamba (State Space Model)
\end{minipage} \\
\midrule\noalign{}
\endhead
\bottomrule\noalign{}
\endlastfoot
\textbf{Asymptotic Complexity} & Quadratic \(O(N^2)\) & \textbf{Linear
\(O(N)\)} \\
\textbf{Indian Pines Train Time / Run} & 75 seconds (avg) & \textbf{26
seconds (avg) --- 2.9× speedup} \\
\textbf{Pavia University Train Time / Run} & 140 seconds (avg) &
\textbf{40 seconds (avg) --- 3.5× speedup} \\
\textbf{Variance / Stability (OA Std)} & \(\pm 3.84\%\) (IP) /
\(\pm 5.98\%\) (PU) & \textbf{\(\pm 1.66\%\) (IP) / \(\pm 0.60\%\) (PU)
--- 10× more stable} \\
\textbf{Convergence Epoch} & Best epoch: 36--48 & \textbf{Best epoch:
27--42 (rapid early convergence)} \\
\end{longtable}

\textbf{Key Finding:} As spatial grid size increased from Indian Pines
(21,025 pixels) to Pavia University (207,400 pixels), MMFormer's runtime
increased by \textbf{\(1.87\times\)}, while its variance inflated to
\(\pm 5.98\%\). In contrast, R2Mamba maintained consistent linear
processing speed and near-zero variance (\(\pm 0.60\%\)).

\begin{center}\rule{0.5\linewidth}{0.5pt}\end{center}

\hypertarget{qualitative-analysis-classification-maps}{%
\subsubsection{4.5 Qualitative Analysis (Classification
Maps)}\label{qualitative-analysis-classification-maps}}

Visual inspection of full-scene prediction maps reveals clear
qualitative distinctions: 1. \textbf{Salt-and-Pepper Artifacts:}
MMFormer produces noticeable isolated misclassifications in uniform crop
fields (e.g., stray soybean pixels inside corn fields). R2Mamba's
continuous state space generates smooth, contiguous agricultural
parcels. 2. \textbf{Edge Boundary Preservation:} Along narrow asphalt
roads and building perimeters in Pavia University, R2Mamba retains sharp
linear fidelity due to its reliability gating, which down-weights
boundary pixel confusion.

\emph{(High-resolution classification maps are archived in
\texttt{results/classification\_maps/})}

\begin{center}\rule{0.5\linewidth}{0.5pt}\end{center}

\hypertarget{chapter-5-conclusion-future-scope}{%
\subsection{Chapter 5: Conclusion \& Future
Scope}\label{chapter-5-conclusion-future-scope}}

\hypertarget{conclusion}{%
\subsubsection{5.1 Conclusion}\label{conclusion}}

This project completed a rigorous empirical benchmarking study of the
two premier 2026 hyperspectral classification architectures:
\textbf{MMFormer} and \textbf{R2Mamba}.

Key conclusions drawn from our experiments: 1. \textbf{Mamba Superiority
in Few-Shot Regimes:} R2Mamba consistently outclassed MMFormer across
both agricultural (Indian Pines: \textbf{+7.48\% OA}) and urban (Pavia
University: \textbf{+5.14\% OA}) scenes. 2. \textbf{Computational
Scalability:} R2Mamba trained \textbf{\(3\times\) to \(3.5\times\)
faster} than MMFormer while consuming substantially less GPU memory due
to its \(O(N)\) linear formulation. 3. \textbf{Statistical Reliability:}
R2Mamba exhibited dramatically lower standard deviation (\(\pm 0.60\%\)
vs \(\pm 5.98\%\)), confirming that selective state transitions are
significantly more robust to small sample fluctuations than multi-head
self-attention.

\hypertarget{future-scope}{%
\subsubsection{5.2 Future Scope}\label{future-scope}}

The modular framework developed in this project provides a foundation
for several promising future extensions: 1. \textbf{Multimodal Fusion:}
Fusing hyperspectral spectral cubes with LiDAR elevation data (e.g.,
Houston18 dataset) to evaluate 3D spatial-elevation reasoning. 2.
\textbf{UAV High-Resolution Mapping:} Benchmarking on ultra-dense
centimeter-scale drone imagery (Wuhan University WHU-Hi series). 3.
\textbf{Planetary Surface Exploration:} Applying trained Mamba models to
NASA MRO CRISM imagery of Mars (Holden and Nili Fossae craters) for
automated mineral identification. 4. \textbf{Edge Deployment:}
Implementing knowledge distillation and FP16/INT8 quantization to deploy
R2Mamba onto embedded satellite edge accelerators.

\begin{center}\rule{0.5\linewidth}{0.5pt}\end{center}

\hypertarget{references}{%
\subsection{References}\label{references}}

\begin{enumerate}
\def\labelenumi{\arabic{enumi}.}
\tightlist
\item
  \textbf{MMFormer:} T. Feng et al., \emph{``Macro-Micro Transformer for
  Hyperspectral Image Classification,''} \emph{IEEE Transactions on
  Geoscience and Remote Sensing (TGRS)}, vol.~64, pp.~1--14, 2026. DOI:
  10.1109/TGRS.2026.3696892.
\item
  \textbf{R2Mamba:} X. Shang et al., \emph{``Route-Reliability Mamba for
  Hyperspectral Image Classification,''} \emph{IEEE Journal of Selected
  Topics in Applied Earth Observations and Remote Sensing (JSTARS)},
  vol.~19, pp.~1--12, 2026. DOI: 10.1109/JSTARS.2026.3728152.
\item
  \textbf{HSI-Models Benchmark Framework:} T. Rachamalla, A. Das, S.
  Kaushik, and S. K. Roy, \emph{``A PyTorch Library for Hyperspectral
  Image Models: Technical Report,''} arXiv preprint arXiv:2609.39871,
  2026.
\item
  \textbf{HybridSN:} G. Roy et al., \emph{``HybridSN: Exploring 3-D--2-D
  CNN Feature Hierarchy for Hyperspectral Image Classification,''}
  \emph{IEEE Geoscience and Remote Sensing Letters (GRSL)}, vol.~17, no.
  2, pp.~277--281, 2019.
\item
  \textbf{SpectralFormer:} D. Hong et al., \emph{``SpectralFormer:
  Rethinking Hyperspectral Image Classification With Transformers,''}
  \emph{IEEE Transactions on Geoscience and Remote Sensing (TGRS)},
  vol.~60, pp.~1--15, 2021.
\item
  \textbf{Mamba (SSM Foundation):} A. Gu and T. Dao, \emph{``Mamba:
  Linear-Time Sequence Modeling with Selective State Spaces,''} arXiv
  preprint arXiv:2312.00752, 2023.
\item
  \textbf{PyTorch:} A. Paszke et al., \emph{``PyTorch: An Imperative
  Style, High-Performance Deep Learning Library,''} in \emph{Advances in
  Neural Information Processing Systems (NeurIPS)}, vol.~32, 2019.
\item
  \textbf{Indian Pines Benchmark:} Purdue University / NASA Jet
  Propulsion Laboratory (JPL) AVIRIS Sensor Dataset.
\item
  \textbf{Pavia University Benchmark:} DLR / University of Pavia ROSIS
  Sensor Optical Benchmark Dataset.
\end{enumerate}

\end{document}
